\chapter{Konstruktion und Gegenbeispiel}
\section{Die Konstruktion}

Die folgende Konstruktion geht auf
\href{https://doi.org/10.1007/BF02389140}%
  {Mogiljanskaja, \emph{Non-isomorphic Semigroups with Isomorphic
  Semigroups of Subsets}, 1973, S.~330--333}
zurück. Ihre Idee ist einfach: Ein zusätzliches Element zerstört die
Isomorphie der Grundhalbgruppen, wird auf der Ebene der Teilmengen aber von
einer unendlichen Reserve aufgenommen. Im Folgenden werden nur die für diesen
Mechanismus wesentlichen Aussagen ausgeführt.

Aus dem vorangehenden Kapitel übernehmen wir die disjunkten Schichten
\[
  L=\{a_i\mid i\in\mathbb N\},
  \qquad
  D=\{b_{ij}\mid i,j\in\mathbb N\}.
\]
Die Eindeutigkeit der \(a\)-Indizes zeigt
\FormulaRefAuto{MogiljanskajaAIndexUnique}[thm].
Für die \(b\)-Indizes gilt \FormulaRefAuto{MogiljanskajaBIndexUnique}[thm].
Die Disjunktheit der Schichten folgt aus
\FormulaRefAuto{MogiljanskajaLayersDisjoint}[thm].

Diese Beweistabellen sind die Referenzquelle für die Konstruktion und
ihre Hilfsaussagen. Die Nummern mit dem Präfix \(28\mathrm E\) gehören zu
dieser Ergänzung. Der Existenzsatz im Hauptwerk steht in
\FormulaRefAuto{MogiljanskajaPairCounterexample}[thm] aus Band~28.
Die \MogReadingLink{} erklärt die Konstruktion und den vollständigen
Beweis in zusammenhängender mathematischer Sprache.

\MogStatementTarget{MogiljanskajaPairDef}
\FormulaDefDeltaK[Die beiden Träger und ihre Multiplikationen]{%
  \begin{gathered}
    \begin{aligned}
      \mathbf0&\coloneqq(2,0),&
      k&\coloneqq(3,0),\\[-2pt]
      S&\coloneqq L\cup D\cup\{\mathbf0\},&
      T&\coloneqq S\cup\{k\}
    \end{aligned}\\[3pt]
    \begin{aligned}[t]
      &\MogStar\colon S\times S\to S,\\[-2pt]
      &\forall i,j\in\mathbb N\,
        \bigl(a_i\MogStar a_j\coloneqq b_{ij}\bigr),\\[-2pt]
      &\forall x,y\in S\,
        \bigl((x\notin L\lor y\notin L)
          \to x\MogStar y\coloneqq\mathbf0\bigr),\\[2pt]
      &\MogDiamond\colon T\times T\to T,\\[-2pt]
      &\forall i,j\in\mathbb N\,
        \bigl(a_i\MogDiamond a_j\coloneqq b_{ij}\bigr),\\[-2pt]
      &\forall x,y\in T\,
        \bigl((x\notin L\lor y\notin L)
          \to x\MogDiamond y\coloneqq\mathbf0\bigr)
    \end{aligned}
  \end{gathered}
}{MogiljanskajaPairDef}{%
  \DeltaRow{Mengen}{L\dsep D\dsep S\dsep T\dsep\mathbf0\dsep k}
  \DeltaRow{Mengen}{a_i\dsep a_j\dsep b_{ij}\dsep i\dsep j\dsep x\dsep y}
  \DeltaRow{Funktionen}{\MogStar\dsep\MogDiamond}
}

\MogProof{MogiljanskajaPairDef}{%
\begin{tabproofwide}
  \proofstepwidestarbreak[]{%
    \MogStar\colon S\times S\to S\dsep
    \MogDiamond\colon T\times T\to T}{\text{\parbox[t]{\linewidth}{\raggedright %
    (M1) Tags, Indizes und Falltrennung}}}
\end{tabproofwide}
}


\MogStatementTarget{MogiljanskajaBasePairProperties}
\FormulaThmDeltaK[Grundlegende Eigenschaften des Mogiljanskaja-Paares]{%
  \begin{aligned}[t]
    &\SemigroupAlg{S}{\MogStar}\land
      \SemigroupAlg{T}{\MogDiamond}\land{}\\[-2pt]
    &\neg\Finite{S}\land\neg\Finite{T}\land{}\\[-2pt]
    &\neg\exists f\,
      \SemigroupIso{f}{S,\MogStar}{T,\MogDiamond}
  \end{aligned}
}{MogiljanskajaBasePairProperties}{%
  \DeltaRow{Mengen}{L\dsep D\dsep S\dsep T\dsep\mathbf0\dsep k}
  \DeltaRow{Mengen}{s\dsep a_i\dsep a_m\dsep a_n\dsep a_0\dsep a_1}
  \DeltaRow{Mengen}{b_{ij}\dsep b_{mn}\dsep b_{i0}\dsep b_{i1}
    \dsep i\dsep m\dsep n}
  \DeltaRow{Funktionen}{a\dsep f\dsep\MogStar\dsep\MogDiamond}
}
\MogProof{MogiljanskajaBasePairProperties}{%
\begin{tabproofwide}
  \MogProofStepWideStarKeep[]{%
    \SemigroupAlg{S}{\MogStar}\land
    \SemigroupAlg{T}{\MogDiamond}}{\text{\parbox[t]{\linewidth}{\raggedright \normalfont (M2): Produktbereich und Nullabsorption}}}
  \MogProofStepWideStarKeep[]{%
    \neg\Finite{S}\land\neg\Finite{T}}{\text{\parbox[t]{\linewidth}{\raggedright \normalfont (M3): injizierte unendliche Schicht}}}
  \MogProofStepWideStarKeep[]{%
    \SemigroupIso{f}{S,\MogStar}{T,\MogDiamond}
      \vdash\exists s\in S\,(f(s)=k)}{\text{\parbox[t]{\linewidth}{\raggedright \normalfont Surjektivität und \ensuremath{k\in T}}}}
  \MogProofStepWideStarKeep[]{%
    \begin{aligned}[t]
      &\SemigroupIso{f}{S,\MogStar}{T,\MogDiamond}
        \dsep s\in S\dsep f(s)=k\dsep s\notin L\\[-2pt]
      &\vdash\exists p\in S\,\exists q\in S\,(s=p\MogStar q)
    \end{aligned}}{\text{\parbox[t]{\linewidth}{\raggedright \normalfont (M4): \ensuremath{s\in D} \normalfont{ oder } \ensuremath{s=\allowbreak \mathbf0}}}}
  \MogProofStepWideStarKeep[4]{%
    \begin{aligned}[t]
      &\SemigroupIso{f}{S,\MogStar}{T,\MogDiamond}
        \dsep s\in S\dsep f(s)=k\dsep s\notin L\\[-2pt]
      &\vdash\exists u\in T\,\exists v\in T\,(k=u\MogDiamond v)
    \end{aligned}}{\text{\parbox[t]{\linewidth}{\raggedright \normalfont Homomorphie; wähle \ensuremath{u=\allowbreak f(p)} \normalfont{ und } \ensuremath{v=\allowbreak f(q)}}}}
  \MogProofStepWideStarKeep[]{%
    \neg\exists u\in T\,\exists v\in T\,(k=u\MogDiamond v)}{\text{\parbox[t]{\linewidth}{\raggedright \normalfont (M4): \ensuremath{k} \normalfont{ liegt nicht im Produktbereich}}}}
  \MogProofStepWideStarKeep[5,6]{%
    \begin{aligned}[t]
      &\SemigroupIso{f}{S,\MogStar}{T,\MogDiamond}
        \dsep s\in S\dsep f(s)=k\\[-2pt]
      &\vdash s\in L
    \end{aligned}}{\text{\parbox[t]{\linewidth}{\raggedright \normalfont Ausschluss der Alternative \ensuremath{s\notin L}}}}
  \MogProofStepWideStarKeep[3,7]{%
    \SemigroupIso{f}{S,\MogStar}{T,\MogDiamond}\vdash
      \exists i\in\mathbb N\,(f(a_i)=k)}{\text{\parbox[t]{\linewidth}{\raggedright \normalfont Surjektivität und Beschreibung von \ensuremath{L}}}}
  \MogProofStepWideStarKeep[8]{%
    \begin{aligned}[t]
      \SemigroupIso{f}{S,\MogStar}{T,\MogDiamond}\vdash
      \exists i\in\mathbb N\,\bigl(&f(b_{i0})=\mathbf0\land{}\\[-2pt]
        &f(b_{i1})=\mathbf0\bigr)
    \end{aligned}}{\text{\parbox[t]{\linewidth}{\raggedright \normalfont (M5): Homomorphie und Multiplikationsvorschrift}}}
  \MogProofStepWideStarKeep[]{i\in\mathbb N\vdash b_{i0}\neq b_{i1}}{\text{\parbox[t]{\linewidth}{\raggedright \normalfont Indexeindeutigkeit und \ensuremath{0\neq1}}}}
  \MogProofStepWideStarKeep[9,10]{%
    \neg\exists f\,\SemigroupIso{f}{S,\MogStar}{T,\MogDiamond}}{\text{\parbox[t]{\linewidth}{\raggedright \normalfont Widerspruch zur Injektivität von \ensuremath{f}}}}
  \MogProofStepWideStarKeep[1,2,11]{%
    \begin{aligned}[t]
      &\SemigroupAlg{S}{\MogStar}\land
        \SemigroupAlg{T}{\MogDiamond}\land{}\\[-2pt]
      &\neg\Finite{S}\land\neg\Finite{T}\land{}\\[-2pt]
      &\neg\exists f\,
        \SemigroupIso{f}{S,\MogStar}{T,\MogDiamond}
    \end{aligned}}{\text{\parbox[t]{\linewidth}{\raggedright \normalfont Bündelung}}}
\end{tabproofwide}
}


\MogStatementTarget{MogiljanskajaProductPartsDef}
\FormulaDefDeltaK[Rechteck- und Nullanteil eines Mengenprodukts]{%
  \begin{aligned}[t]
    R(X,Y)&\coloneqq
      \bigl\{b_{ij}\bigm|
        i,j\in\mathbb N,\,
        a_i\in X\cap L,\,
        a_j\in Y\cap L\bigr\},\\[-2pt]
    \varepsilon(X,Y)&\coloneqq
      \begin{cases}
        \varnothing,&X\subseteq L\land Y\subseteq L,\\
        \{\mathbf0\},&\text{sonst}
      \end{cases}
  \end{aligned}
}{MogiljanskajaProductPartsDef}{%
  \DeltaRow{Mengen}{L\dsep D\dsep X\dsep Y\dsep
    a_i\dsep a_j\dsep b_{ij}\dsep i\dsep j
    \dsep R(X,Y)\dsep\varepsilon(X,Y)}
}

\MogStatementTarget{MogiljanskajaPowerProductFormula}
\FormulaThmDeltaKR[Direkte Formel für die beiden Mengenprodukte]{%
 \begin{aligned}[t]
  &\text{(i)}\quad X,Y\in\PowNE{S} \vdash X\MogStar_{\mathcal P}Y =R(X,Y)\cup\varepsilon(X,Y)\\[-2pt]
  &\text{(ii)}\quad X,Y\in\PowNE{T} \vdash X\MogDiamond_{\mathcal P}Y =R(X,Y)\cup\varepsilon(X,Y)
 \end{aligned}%
}{%
  \begin{aligned}[t]
    X,Y\in\PowNE{S}
      &\vdash X\MogStar_{\mathcal P}Y
        =R(X,Y)\cup\varepsilon(X,Y),\\[-2pt]
    X,Y\in\PowNE{T}
      &\vdash X\MogDiamond_{\mathcal P}Y
        =R(X,Y)\cup\varepsilon(X,Y)
  \end{aligned}
}{MogiljanskajaPowerProductFormula}{%
  \DeltaRow{Mengen}{L\dsep D\dsep S\dsep T\dsep\mathbf0
    \dsep X\dsep Y\dsep R(X,Y)\dsep\varepsilon(X,Y)}
  \DeltaRow{Funktionen}{\MogStar_{\mathcal P}
    \dsep\MogDiamond_{\mathcal P}}
}
\MogProof{MogiljanskajaPowerProductFormula}{%
\begin{tabproofwide}
  \proofstepwidestarbreak[]{%
    \begin{aligned}[t]
      &X,Y\in\PowNE{S}\vdash{}\\[-2pt]
      &\quad X\MogStar_{\mathcal P}Y
          =R(X,Y)\cup\varepsilon(X,Y),\\[-2pt]
      &X,Y\in\PowNE{T}\vdash{}\\[-2pt]
      &\quad X\MogDiamond_{\mathcal P}Y
          =R(X,Y)\cup\varepsilon(X,Y)
    \end{aligned}}{%
    \text{\parbox[t]{\linewidth}{\raggedright \normalfont(M6) Paarweise Produktklassifikation}}}
\end{tabproofwide}
}


\Needspace{10\baselineskip}
Für die Reserve verwenden wir die im vorangehenden Kapitel definierten Mengen und
Bijektionen
\[
  \begin{gathered}
    c_0=b_{00},\qquad c_1=b_{11},\qquad c_2=b_{10},\\
    U=\{b_{(n+1)+1,j}\mid n,j\in\mathbb N\},\qquad
    V=U\cup\{c_2\},\\
    P=\{c_0,c_1\},\qquad D'=P\cup V,\\
    \sigma\colon U\bij V,\qquad
    \theta\colon U\bij D.
  \end{gathered}
\]
Dies sind \FormulaRefAuto{MogiljanskajaDPrimeDef}[def],
\FormulaRefAuto{MogiljanskajaSigmaBijective}[thm] und
\FormulaRefAuto{MogiljanskajaThetaBijective}[thm]. Nach
\FormulaRefAuto{MogiljanskajaReserveMarkerFacts}[thm] und
\FormulaRefAuto{MogiljanskajaReserveCoreFacts}[thm] gilt insbesondere
\[
  P\cap U=\varnothing,\qquad
  c_2\notin P\cup U,\qquad
  D'\subseteq D,\qquad
  b_{01}\notin D'.
\]

\MogStatementTarget{MogiljanskajaReserveMapsDef}
\FormulaDefDeltaK[Reserve und lokale Potenzmengenabbildung]{%
  \begin{aligned}[t]
    \mathcal C&\coloneqq\CoreFamily{P}{V},&
    \mathcal C_0&\coloneqq\CoreFamily{P}{U},\\[-2pt]
    \mathcal C_1&\coloneqq\CoreFamily{P\cup\{c_2\}}{U},&
    \mathcal E&\coloneqq\CoreFamily{\{k\}}{D},\\[-2pt]
    \kappa_0&\coloneqq\CoreTransport{P}{P}{\sigma},&
    \kappa_1&\coloneqq
      \CoreTransport{P\cup\{c_2\}}{\{k\}}{\theta},\\[-2pt]
    \psi&\coloneqq
      \CaseFun{\mathcal C_0}{\mathcal C_1}{\kappa_0}{\kappa_1},\\[-2pt]
    \mathcal O&\coloneqq\powerset(D)\setminus\mathcal C,&
    \eta&\coloneqq
      \CaseFun{\mathcal O}{\mathcal C}{\Id_{\mathcal O}}{\psi}
  \end{aligned}
}{MogiljanskajaReserveMapsDef}{%
  \DeltaRow{Mengen}{D\dsep U\dsep V\dsep P\dsep c_2\dsep k
    \dsep\mathcal C\dsep\mathcal C_0\dsep\mathcal C_1
    \dsep\mathcal E\dsep\mathcal O}
  \DeltaRow{Funktionen}{\sigma\dsep\theta\dsep
    \kappa_0\dsep\kappa_1\dsep\psi\dsep\eta}
}

\MogStatementTarget{MogiljanskajaReserveMapProperties}
\FormulaThmDeltaKR[Eigenschaften der Reserveabbildung]{%
 \begin{aligned}[t]
  &\text{(i)}\quad \psi\colon\mathcal C\bij\mathcal C\cup\mathcal E\\[-2pt]
  &\text{(ii)}\quad \eta\colon\powerset(D)\bij\powerset(D\cup\{k\})\\[-2pt]
  &\text{(iii)}\quad \eta(\varnothing)=\varnothing\\[-2pt]
  &\text{(iv)}\quad K\in\powerset(D)\dsep K\notin\mathcal C \vdash\eta(K)=K\\[-2pt]
  &\text{(v)}\quad X,Y\in\powerset(T)\vdash R(X,Y)\notin\mathcal C
 \end{aligned}%
}{%
  \begin{aligned}[t]
    &\psi\colon\mathcal C\bij\mathcal C\cup\mathcal E,\\[-2pt]
    &\eta\colon\powerset(D)\bij\powerset(D\cup\{k\})
      \dsep\eta(\varnothing)=\varnothing,\\[-2pt]
    &K\in\powerset(D)\dsep K\notin\mathcal C
      \vdash\eta(K)=K,\\[-2pt]
    &X,Y\in\powerset(T)\vdash R(X,Y)\notin\mathcal C
  \end{aligned}
}{MogiljanskajaReserveMapProperties}{%
  \DeltaRow{Mengen}{D\dsep D'\dsep P\dsep U\dsep V\dsep k
    \dsep\mathcal C\dsep\mathcal E\dsep\mathcal O
    \dsep K\dsep X\dsep Y\dsep R(X,Y)}
  \DeltaRow{Funktionen}{\psi\dsep\eta}
}
\MogProof{MogiljanskajaReserveMapProperties}{%
\begin{tabproofwide}
  \proofstepwidestar[]{%
    \psi\colon\mathcal C\bij\mathcal C\cup\mathcal E}{\text{\parbox[t]{\linewidth}{\raggedright %
    (M7)--(M8) Kerntransport und Disjunktheit}}}
  \proofstepwidestar[]{%
    \eta\colon\powerset(D)\bij\powerset(D\cup\{k\})}{\text{\parbox[t]{\linewidth}{\raggedright %
    (M9) disjunkte Zielzerlegung}}}
  \proofstepwidestar[]{%
    \eta(\varnothing)=\varnothing}{\text{\parbox[t]{\linewidth}{\raggedright %
    (M10) geschützter unveränderter Teil}}}
  \proofstepwidestar[]{%
    K\in\powerset(D)\dsep K\notin\mathcal C
      \vdash\eta(K)=K}{\text{\parbox[t]{\linewidth}{\raggedright %
    (M10) Identitätszweig}}}
  \proofstepwidestar[]{%
    X,Y\in\powerset(T)\vdash R(X,Y)\notin\mathcal C}{\text{\parbox[t]{\linewidth}{\raggedright %
    (M11) erzwungene Ecke \(b_{01}\)}}}
\end{tabproofwide}
}


\MogStatementTarget{MogiljanskajaPhiDef}
\FormulaDefDeltaK[Die globale Potenzmengenabbildung]{%
  \begin{aligned}[t]
    A_*&\coloneqq L\cup\{\mathbf0\},\\[-2pt]
    \Phi&\coloneqq
      \LayerPowerMap{A_*}{D}{D\cup\{k\}}{\eta}
      \restriction_{\PowNE{S}}
  \end{aligned}
}{MogiljanskajaPhiDef}{%
  \DeltaRow{Mengen}{A_*\dsep L\dsep D\dsep S\dsep T
    \dsep\mathbf0\dsep k}
  \DeltaRow{Funktionen}{\eta\dsep\Phi}
}

Für \(X\in\PowNE{S}\) lautet die Grundgleichung
\[
  \Phi(X)=(X\cap A_*)\cup\eta(X\cap D).
\]

\MogStatementTarget{MogiljanskajaPowerSemigroupIsomorphism}
\FormulaThmDeltaK[Isomorphie der Potenzhalbgruppen]{%
  \SemigroupIso{\Phi}
    {\PowNE{S},\MogStar_{\mathcal P}}
    {\PowNE{T},\MogDiamond_{\mathcal P}}%
}{MogiljanskajaPowerSemigroupIsomorphism}{%
  \DeltaRow{Mengen}{A_*\dsep L\dsep D\dsep S\dsep T
    \dsep X\dsep Y\dsep Z}
  \DeltaRow{Funktionen}{\eta\dsep\Phi
    \dsep\MogStar_{\mathcal P}\dsep\MogDiamond_{\mathcal P}}
}
\MogProof{MogiljanskajaPowerSemigroupIsomorphism}{%
\begin{tabproofwide}
  \proofstepwidestar[]{%
    \Phi\colon\PowNE{S}\bij\PowNE{T}}{\text{\parbox[t]{\linewidth}{\raggedright %
    (M12) nulltreue Schichtfortsetzung}}}
  \proofstepwidestarbreak[]{%
    \begin{aligned}[t]
      \Phi(X)\cap L&=X\cap L,\\[-2pt]
      X\subseteq L&\Longleftrightarrow\Phi(X)\subseteq L
    \end{aligned}}{%
    \text{\parbox[t]{\linewidth}{\raggedright \normalfont(M13) Invarianz der geschützten Schicht}}}
  \proofstepwidestarbreak[]{%
    \Phi(X)\MogDiamond_{\mathcal P}\Phi(Y)
      =X\MogStar_{\mathcal P}Y}{\text{\parbox[t]{\linewidth}{\raggedright \normalfont(M14) Invarianz von \ensuremath{R} \normalfont{ und } \ensuremath{\varepsilon}}}}
  \proofstepwidestarbreak[]{%
    \Phi(X\MogStar_{\mathcal P}Y)
      =X\MogStar_{\mathcal P}Y}{%
    \text{\parbox[t]{\linewidth}{\raggedright \normalfont(M15) Produktrechtecke sind geschützt}}}
  \proofstepwidestarbreak[3,4]{%
    \Phi(X\MogStar_{\mathcal P}Y)
      =\Phi(X)\MogDiamond_{\mathcal P}\Phi(Y)}{%
    \text{\parbox[t]{\linewidth}{\raggedright \normalfont(M16) Vergleich der Fixpunktgleichungen}}}
  \proofstepwidestarbreak[1,5]{%
    \SemigroupIso{\Phi}
      {\PowNE{S},\MogStar_{\mathcal P}}
      {\PowNE{T},\MogDiamond_{\mathcal P}}}{%
    \text{\parbox[t]{\linewidth}{\raggedright \normalfont(M16) Bijektivität und Homomorphie}}}
\end{tabproofwide}
}


\MogStatementTarget{MogiljanskajaConstructedPairCounterexample}
\FormulaThmDeltaK[Das Mogiljanskaja-Paar ist ein unendliches
Gegenbeispiel]{%
  \begin{aligned}[t]
    &\SemigroupAlg{S}{\MogStar}\land
      \SemigroupAlg{T}{\MogDiamond}\land{}\\[-2pt]
    &\neg\Finite{S}\land\neg\Finite{T}\land{}\\[-2pt]
    &\neg\exists f\,
      \SemigroupIso{f}{S,\MogStar}{T,\MogDiamond}\land{}\\[-2pt]
    &\exists F\,\SemigroupIso{F}
      {\PowNE{S},\MogStar_{\mathcal P}}
      {\PowNE{T},\MogDiamond_{\mathcal P}}
  \end{aligned}
}{MogiljanskajaConstructedPairCounterexample}{%
  \DeltaRow{Mengen}{S\dsep T}
  \DeltaRow{Funktionen}{f\dsep F\dsep\Phi
    \dsep\MogStar\dsep\MogDiamond
    \dsep\MogStar_{\mathcal P}\dsep\MogDiamond_{\mathcal P}}
}
\MogProof{MogiljanskajaConstructedPairCounterexample}{%
\begin{tabproofwide}
  \proofstepwidestar[]{%
    \begin{aligned}[t]
      &\SemigroupAlg{S}{\MogStar}\land
        \SemigroupAlg{T}{\MogDiamond}\land{}\\[-2pt]
      &\neg\Finite{S}\land\neg\Finite{T}\land{}\\[-2pt]
      &\neg\exists f\,
        \SemigroupIso{f}{S,\MogStar}{T,\MogDiamond}
    \end{aligned}}{%
    \FormulaRefAuto{MogiljanskajaBasePairProperties}[thm]}
  \proofstepwidestarbreak[]{%
    \exists F\,\SemigroupIso{F}
      {\PowNE{S},\MogStar_{\mathcal P}}
      {\PowNE{T},\MogDiamond_{\mathcal P}}}{\text{\parbox[t]{\linewidth}{\raggedright \normalfont(M17) Zeuge \ensuremath{F=\allowbreak \Phi}}}}
  \proofstepwidestar[1,2]{%
    \begin{aligned}[t]
      &\SemigroupAlg{S}{\MogStar}\land
        \SemigroupAlg{T}{\MogDiamond}\land{}\\[-2pt]
      &\neg\Finite{S}\land\neg\Finite{T}\land{}\\[-2pt]
      &\neg\exists f\,
        \SemigroupIso{f}{S,\MogStar}{T,\MogDiamond}\land{}\\[-2pt]
      &\exists F\,\SemigroupIso{F}
        {\PowNE{S},\MogStar_{\mathcal P}}
        {\PowNE{T},\MogDiamond_{\mathcal P}}
    \end{aligned}}{\text{\parbox[t]{\linewidth}{\raggedright \normalfont Bündelung}}}
\end{tabproofwide}
}


Die konstruierten Träger \(S,T\) und Operationen
\(\MogStar,\MogDiamond\) liefern damit ausdrücklich die Zeugen für den
Existenzsatz \FormulaRefAuto{MogiljanskajaPairCounterexample}[thm]
in Band~28: Die soeben bewiesene Konjunktion erfüllt sämtliche dort
verlangten Eigenschaften. Viermalige Existenz-Einführung ergibt die
Behauptung über die Existenz von \(S,T,\MogStar,\MogDiamond\).


Damit ist das Mogiljanskaja-Paar ein Gegenbeispiel zur Umkehrung von
\FormulaRefAuto{SemigroupIsoImpliesPowerSemigroupIso}[thm].
Die Umkehrung ist also ohne Endlichkeitsvoraussetzung falsch.

